\documentclass[11pt,a4paper]{article}

% --- Encoding ---
\usepackage[utf8]{inputenc}
\usepackage[T1]{fontenc}
\usepackage[spanish,es-noshorthands]{babel}
\usepackage{lmodern}

% --- Layout ---
\usepackage[margin=2.5cm]{geometry}
\usepackage{parskip}
\setlength{\parindent}{0pt}
\setlength{\parskip}{0.5em plus 0.2em minus 0.1em}

% --- Tables and graphics ---
\usepackage{booktabs}
\usepackage{array}
\usepackage{enumitem}
\usepackage{xcolor}
\usepackage{graphicx}
\usepackage{tikz}
\usetikzlibrary{positioning,shapes.geometric,arrows.meta,fit,backgrounds,calc,decorations.pathreplacing}
\usepackage{caption}
\usepackage{subcaption}
\usepackage{amsmath,amssymb}

% --- Soporte para caracteres Unicode usados en formulas en texto plano ---
\usepackage{newunicodechar}
\newunicodechar{→}{\ensuremath{\rightarrow}}
\newunicodechar{≥}{\ensuremath{\geq}}
\newunicodechar{≤}{\ensuremath{\leq}}
\newunicodechar{×}{\ensuremath{\times}}
\newunicodechar{÷}{\ensuremath{\div}}
\newunicodechar{≈}{\ensuremath{\approx}}
\newunicodechar{⌈}{\ensuremath{\lceil}}
\newunicodechar{⌉}{\ensuremath{\rceil}}
\newunicodechar{✓}{\ensuremath{\checkmark}}
\newunicodechar{✗}{\ensuremath{\times}}

% --- Code listings ---
\usepackage{listings}
\definecolor{codegray}{rgb}{0.96,0.96,0.96}
\definecolor{codeborder}{rgb}{0.85,0.85,0.85}
\definecolor{codecomment}{rgb}{0.30,0.55,0.30}
\definecolor{codekeyword}{rgb}{0.10,0.30,0.70}
\definecolor{codestring}{rgb}{0.65,0.15,0.10}

\lstdefinestyle{cstyle}{
    language=C,
    backgroundcolor=\color{codegray},
    basicstyle=\ttfamily\footnotesize,
    breaklines=true,
    frame=single,
    rulecolor=\color{codeborder},
    showstringspaces=false,
    keywordstyle=\color{codekeyword}\bfseries,
    commentstyle=\color{codecomment}\itshape,
    stringstyle=\color{codestring},
    aboveskip=1em,
    belowskip=1em,
}

\lstdefinestyle{shellstyle}{
    language=bash,
    backgroundcolor=\color{codegray},
    basicstyle=\ttfamily\small,
    breaklines=true,
    frame=single,
    rulecolor=\color{codeborder},
    showstringspaces=false,
    keywordstyle=\color{codekeyword}\bfseries,
    commentstyle=\color{codecomment}\itshape,
    stringstyle=\color{codestring},
    aboveskip=1em,
    belowskip=1em,
    columns=fullflexible,
    upquote=true,
}

\lstdefinestyle{textstyle}{
    backgroundcolor=\color{codegray},
    basicstyle=\ttfamily\footnotesize,
    breaklines=true,
    frame=single,
    rulecolor=\color{codeborder},
    showstringspaces=false,
    aboveskip=1em,
    belowskip=1em,
    columns=fullflexible,
}

\lstset{style=textstyle}

% --- Definition / Idea / Important boxes ---
\usepackage[most]{tcolorbox}
\newtcolorbox{idea}[1][]{
    colback=orange!5!white,
    colframe=orange!60!black,
    fonttitle=\bfseries,
    title=Idea intuitiva,
    sharp corners=southwest,
    rounded corners=northwest,
    arc=2pt,
    boxrule=0.6pt,
    left=8pt,right=8pt,top=4pt,bottom=4pt,
    #1
}
\newtcolorbox{importante}[1][]{
    colback=red!5!white,
    colframe=red!50!black,
    fonttitle=\bfseries,
    title=Importante,
    sharp corners=southwest,
    rounded corners=northwest,
    arc=2pt,
    boxrule=0.6pt,
    left=8pt,right=8pt,top=4pt,bottom=4pt,
    #1
}
\newtcolorbox{respuesta}[1][]{
    colback=green!5!white,
    colframe=green!40!black,
    fonttitle=\bfseries,
    title=Respuesta,
    sharp corners=southwest,
    rounded corners=northwest,
    arc=2pt,
    boxrule=0.6pt,
    left=8pt,right=8pt,top=4pt,bottom=4pt,
    #1
}

% --- Hyperlinks ---
\usepackage[colorlinks=true,
            linkcolor=blue!60!black,
            urlcolor=blue!50!black,
            citecolor=blue!60!black]{hyperref}

% --- Color palette ---
\definecolor{c1}{HTML}{4C72B0}
\definecolor{c2}{HTML}{DD8452}
\definecolor{c3}{HTML}{55A868}
\definecolor{c4}{HTML}{8172B2}
\definecolor{cred}{HTML}{C44E52}

% --- Title block ---
\title{
    \vspace{-2.5em}
    \textbf{Parcial 1 --- 2026} \\[0.3em]
    \large Sistemas Operativos / Sistemas Operativos y Redes
}
\author{
    Tomás Rodeghiero \\
    \small UNRC --- Universidad Nacional de Río Cuarto
}
\date{2026}

\begin{document}

\maketitle

\begin{abstract}
\noindent
Resolución completa del Parcial 1 del año 2026 de la cátedra de
Sistemas Operativos (Lic.\ en Ciencias de la Computación, UNRC). El
parcial consta de cinco ejercicios que abarcan: (i) verdadero/falso
sobre arquitectura de SO (monolítico vs.\ microkernel, traps, modo
usuario/kernel, multitarea, procesos zombie, shell); (ii) linker
estático y resolución de símbolos entre dos archivos objeto; (iii)
posibles salidas de un programa con \texttt{fork()}/\texttt{wait()}
contrastado con uno multihilo que comparte una variable; (iv) análisis
de \emph{deadlock} con grafo de dependencias entre tres tareas y dos
recursos; (v) planificación basada en prioridades con \emph{preemption}
y \emph{quantum} sobre tres tareas. Cada ejercicio se desarrolla con
estructura uniforme: \textbf{Enunciado $\to$ Análisis $\to$ Desarrollo
$\to$ Respuesta}.
\end{abstract}

\tableofcontents
\bigskip

% =====================================================================
\section{Ejercicio 1 --- Verdadero o falso (2 pts)}
% =====================================================================

\textbf{Enunciado.} Marcar verdad (V) o falsedad (F) de las siguientes
sentencias. Justificar brevemente las falsas.

\subsection*{1.a --- SO monolítico}

\emph{``Un SO monolítico tiene sus subsistemas o módulos en un único
programa. La comunicación entre módulos se hace mediante variables
globales e invocaciones a funciones.''}

\begin{respuesta}
\textbf{V}. En un núcleo monolítico (como Linux clásico) todos los
subsistemas ---scheduler, gestor de memoria, sistema de archivos,
drivers, networking--- se compilan en un único binario y corren en el
mismo espacio de direcciones, en modo kernel. La comunicación entre
ellos es directa: invocación de funciones y acceso a estructuras
globales (tabla de procesos, tabla de archivos abiertos, listas de
páginas libres, etc.).
\end{respuesta}

\subsection*{1.b --- Microkernel y rendimiento}

\emph{``Un SO microkernel logra mejor rendimiento sobre uno monolítico
debido a su kernel minimalista.''}

\begin{respuesta}
\textbf{F}. Un microkernel típicamente tiene \emph{peor} rendimiento
que un monolítico, no mejor. En un microkernel los subsistemas
(filesystem, drivers, gestor de memoria) corren como procesos
servidores en \textbf{modo usuario} y se comunican entre sí por
\emph{message passing} (IPC). Cada operación que en un monolítico es
una llamada a función dentro del mismo espacio de direcciones, en un
microkernel implica al menos dos cambios de modo, dos cambios de
contexto y la copia de los datos entre espacios.

Las ventajas del microkernel son \textbf{modularidad}, \textbf{robustez}
(un driver que crashea no tira al kernel) y \textbf{seguridad}, pero
\textbf{no} performance.
\end{respuesta}

\subsection*{1.c --- La instrucción de \texttt{trap}}

\emph{``La instrucción de \texttt{trap} (usada para hacer un syscall)
debe ser una instrucción privilegiada.''}

\begin{respuesta}
\textbf{F}. Justamente lo contrario: la instrucción de \texttt{trap}
(\texttt{syscall} en x86-64, \texttt{ecall} en RISC-V, \texttt{svc} en
ARM) \emph{no} es privilegiada. Si lo fuera, ningún programa de
usuario podría usarla, lo cual rompería todo el modelo de syscalls.

La instrucción de trap es una de las pocas formas legítimas en que
un proceso de usuario puede provocar un cambio a modo supervisor.
Hardware mediante: la CPU detecta la instrucción, salva el estado,
cambia a modo kernel y salta a una dirección fija (la \emph{trap
table}) donde el kernel toma el control. Lo \emph{privilegiado} es lo
que ejecuta el kernel después del trap; la instrucción en sí está
diseñada para ser disponible desde modo usuario.
\end{respuesta}

\subsection*{1.d --- Timesharing y \emph{quantums}}

\emph{``Un SO \emph{timesharing} multiplexa la CPU entre varias tareas
con usos basados en \emph{quantums}.''}

\begin{respuesta}
\textbf{V}. Un SO de tiempo compartido asigna a cada tarea un cuanto
de tiempo de CPU (\emph{quantum}). Al vencer el cuanto, el timer
dispara una interrupción que invoca al scheduler para decidir qué
tarea ejecutar a continuación. Este es exactamente el modelo
\emph{preemptive} sobre el que se construyen Round Robin y todos los
planificadores modernos de tiempo compartido.
\end{respuesta}

\subsection*{1.e --- Multitarea y CPUs}

\emph{``Un SO \emph{multitarea} sólo es posible con múltiples CPUs.''}

\begin{respuesta}
\textbf{F}. Multitarea \textbf{no} requiere múltiples CPUs. Se puede
implementar perfectamente en \textbf{una sola CPU} mediante
multiplexación temporal (timesharing). De hecho, fue así por décadas
antes de la era multicore: Linux, UNIX y Windows corrían multitarea en
un único procesador.

La distinción clave es:

\begin{itemize}[leftmargin=*,nosep]
    \item \textbf{Multitarea / concurrencia}: \emph{apariencia} de
        ejecución simultánea de varias tareas. Se logra con una sola
        CPU multiplexando entre ellas.
    \item \textbf{Paralelismo real}: ejecución \emph{verdaderamente
        simultánea} de varias tareas. Esto sí requiere múltiples CPUs
        (o cores).
\end{itemize}
\end{respuesta}

\subsection*{1.f --- Proceso \emph{zombie}}

\emph{``Un proceso Zombie ya finalizó su ejecución.''}

\begin{respuesta}
\textbf{V}. Un proceso zombie es uno que ya terminó (ejecutó
\texttt{exit()} o fue terminado por una señal) pero cuyo descriptor
(PCB) sigue existiendo en la tabla de procesos del kernel porque su
\emph{padre} aún no llamó a \texttt{wait()}/\texttt{waitpid()} para
recoger su exit status. El proceso ``está muerto'' (no consume CPU, no
tiene memoria asignada salvo el descriptor) pero su entrada persiste.
Se ve en \texttt{ps} con estado \texttt{Z}.
\end{respuesta}

\subsection*{1.g --- Acceso de un proceso de usuario al reloj}

\emph{``Un proceso de usuario puede acceder al dispositivo del reloj y
programarlo para que genere una interrupción luego de un intervalo de
tiempo dado.''}

\begin{respuesta}
\textbf{F}. Un proceso de usuario \textbf{no} puede acceder
directamente al hardware del reloj. Los registros del timer y de las
IRQs están en el espacio de direcciones del kernel y solo son
accesibles en modo supervisor. La MMU bloquea cualquier intento de
lectura/escritura directa desde modo usuario.

Lo que el proceso de usuario sí puede hacer es \textbf{pedírselo al
kernel} vía syscalls: \texttt{alarm(N)}, \texttt{setitimer},
\texttt{timer\_create}, etc. El kernel programa el hardware en su
nombre y, cuando vence el timer, envía una señal (\texttt{SIGALRM},
por ejemplo) al proceso. La interacción es siempre \emph{indirecta},
mediada por el kernel.
\end{respuesta}

\subsection*{1.h --- El shell}

\emph{``El \emph{shell} es un módulo del kernel.''}

\begin{respuesta}
\textbf{F}. El shell (\texttt{bash}, \texttt{zsh}, \texttt{fish}) es
un \textbf{programa de usuario} normal y corriente, no parte del
kernel. Es simplemente un proceso que corre en modo usuario igual que
cualquier otra aplicación: lee comandos por \texttt{stdin}, los
interpreta, hace \texttt{fork()}/\texttt{exec()} para lanzar otros
procesos, y maneja redirecciones, pipes y señales \emph{usando
syscalls como cualquier otro programa}. No tiene ningún privilegio
especial salvo los del usuario que lo lanzó.
\end{respuesta}

% =====================================================================
\section{Ejercicio 2 --- Linking estático (2 pts)}
% =====================================================================

\textbf{Enunciado.} Dados los dos archivos objeto siguientes,
enlazados estáticamente con \texttt{gcc a.o b.o -o myprog}, determinar
si las afirmaciones son verdaderas o falsas.

\subsection*{Contenido de los archivos objeto}

\begin{center}
\begin{minipage}[t]{0.46\textwidth}
\textbf{a.o}
\begin{lstlisting}[style=textstyle]
.text
  0x0000  ...
  0x00A0  call 0x00A0
  0x00A4  ...
  ...
  // end .text section
.data
  0x00E0  ...
.relocation entries
  addr     type   symbol
  0x00A0   call   f
\end{lstlisting}
\end{minipage}
\hfill
\begin{minipage}[t]{0.46\textwidth}
\textbf{b.o}
\begin{lstlisting}[style=textstyle]
.text
  0x0000  ...
  0x001A  f: ...
  ...
  0x00A0  ret
  // end .text section
.data
  0x00A4  ...
.symbol table
  symbol   address
  f        0x001A
\end{lstlisting}
\end{minipage}
\end{center}

\subsection*{Análisis general}

Antes de resolver los tres ítems, recordemos qué hace el linker en un
linking estático:

\begin{enumerate}[leftmargin=*,nosep]
    \item \textbf{Concatena} las secciones del mismo nombre de todos
        los archivos objeto. Las \texttt{.text} de \texttt{a.o} y
        \texttt{b.o} forman, una a continuación de la otra, la
        \texttt{.text} de \texttt{myprog}.
    \item \textbf{Asigna} direcciones absolutas a cada sección
        concatenada en el ejecutable final.
    \item \textbf{Resuelve} las referencias externas: para cada
        entrada en la tabla de relocaciones, ubica el símbolo en la
        tabla de símbolos del otro \texttt{.o} y completa la
        instrucción con la dirección final.
\end{enumerate}

\textbf{Mapa de direcciones en \texttt{myprog}.} Tomando la
\texttt{.text} de \texttt{a.o} primero (tamaño \texttt{0x00E0}, ya que
\texttt{.data} comienza ahí en \texttt{a.o}), y luego la
\texttt{.text} de \texttt{b.o} (tamaño \texttt{0x00A4}):

\begin{center}
\begin{tikzpicture}[every node/.style={font=\ttfamily\small},
    sec/.style={draw,minimum height=0.75cm,align=center}]
\node[sec,fill=c1!20,minimum width=4.0cm] (a) at (0,0) {.text de a.o};
\node[sec,fill=c2!20,minimum width=3.2cm,right=0pt of a] (b) {.text de b.o};
\node[font=\scriptsize] at ($(a.south west)+(0.1,-0.25)$) {0x0000};
\node[font=\scriptsize] at ($(a.south east)+(0,-0.25)$) {0x00E0};
\node[font=\scriptsize] at ($(b.south east)+(0,-0.25)$) {0x0184};
\node[font=\scriptsize] at ($(a.north)+(0,0.25)$) {0xE0 bytes};
\node[font=\scriptsize] at ($(b.north)+(0,0.25)$) {0xA4 bytes};
\end{tikzpicture}
\end{center}

Es decir, en \texttt{myprog}:

\begin{itemize}[leftmargin=*,nosep]
    \item La \texttt{.text} de \texttt{a.o} ocupa
        $[\texttt{0x0000}, \texttt{0x00E0})$.
    \item La \texttt{.text} de \texttt{b.o} ocupa
        $[\texttt{0x00E0}, \texttt{0x0184})$.
\end{itemize}

La función \texttt{f}, que en \texttt{b.o} estaba en el offset
\texttt{0x001A} dentro de su propia \texttt{.text}, queda relocada a
$\texttt{0x00E0} + \texttt{0x001A} = \texttt{0x00FA}$ en
\texttt{myprog}.

\subsection*{2.a --- Dirección de \texttt{f} en \texttt{myprog}}

\emph{``La dirección de la función \texttt{f} en \texttt{myprog}
resultará \texttt{0x001A}.''}

\begin{respuesta}
\textbf{F}. \texttt{0x001A} es solo el offset de \texttt{f} dentro de
la \texttt{.text} de \texttt{b.o}; es una dirección \emph{relativa} en
el archivo objeto antes del linking. Al hacer
\texttt{gcc a.o b.o -o myprog}, el linker coloca la \texttt{.text} de
\texttt{b.o} a continuación de la de \texttt{a.o} (que mide
\texttt{0xE0}). La dirección absoluta de \texttt{f} en el ejecutable
final es:

\quad dir(\texttt{f}) en myprog = inicio de b.o.text + offset interno
                                 = \texttt{0x00E0} + \texttt{0x001A}
                                 = \textbf{\texttt{0x00FA}}.
\end{respuesta}

\subsection*{2.b --- ¿\texttt{call 0x00A0} es una llamada local?}

\emph{``La instrucción \texttt{call 0x00A0} es una llamada a una
función local en \texttt{a.o}.''}

\begin{respuesta}
\textbf{F}. No es una llamada local. La pista definitiva está en la
\textbf{tabla de relocaciones} de \texttt{a.o}:

\begin{lstlisting}[style=textstyle]
.relocation entries
  addr     type   symbol
  0x00A0   call   f
\end{lstlisting}

Esta entrada dice: ``en la dirección \texttt{0x00A0} hay una
instrucción de tipo \texttt{call} cuyo operando debe completarse con
la dirección final del símbolo externo \texttt{f}''. El \texttt{0x00A0}
que aparece como operando del call \emph{es un \texttt{placeholder}}
(0 o autorreferencia) que se rellena durante el linking. El símbolo
\texttt{f} \textbf{no} está definido en \texttt{a.o} sino en
\texttt{b.o}, por lo cual la llamada es \emph{externa}, no local.

Si fuera local, la referencia ya estaría resuelta en \texttt{a.o} y
no aparecería en la tabla de relocaciones.
\end{respuesta}

\subsection*{2.c --- Tamaño de \texttt{.text} en \texttt{myprog}}

\emph{``La sección \texttt{.text} de \texttt{myprog} será de tamaño
\texttt{0x00E0 + 0x00A4}.''}

\begin{respuesta}
\textbf{V} (con la salvedad de alineación que se aclara abajo).

En \texttt{a.o}, la sección \texttt{.text} arranca en \texttt{0x0000}
y la sección \texttt{.data} comienza inmediatamente en \texttt{0x00E0}.
Es decir, el \emph{tamaño} de \texttt{.text} en \texttt{a.o} es
\texttt{0xE0} bytes. Análogamente, en \texttt{b.o} la \texttt{.text}
mide \texttt{0xA4} bytes (desde \texttt{0x0000} hasta donde comienza
\texttt{.data} en \texttt{0x00A4}).

Como el linker concatena las dos \texttt{.text} en una sola, el
tamaño resultante es la \textbf{suma}:

\quad \texttt{0x00E0} + \texttt{0x00A4} = \texttt{0x0184} bytes.

\textbf{Caveat real}: en sistemas reales el linker puede insertar
algunos bytes de \emph{padding} para satisfacer requisitos de
alineación entre las dos secciones \texttt{.text} concatenadas (en
x86 típicamente 16 bytes), por lo que el tamaño \emph{exacto} podría
ser un poco mayor. Pero la suma es esencialmente correcta y la
afirmación es válida en el modelo simplificado del ejercicio.
\end{respuesta}

% =====================================================================
\section{Ejercicio 3 --- Salidas posibles: \texttt{fork} vs.\ threads (2 pts)}
% =====================================================================

\textbf{Enunciado.} Mostrar las posibles salidas de los dos programas.

\subsection*{Programa A (con \texttt{fork})}

\begin{lstlisting}[style=cstyle]
int x = 0;
if (fork() == 0) {
    /* hijo */
    x += 1;
    printf("child:%d ", x);
    exit(x);
} else {
    /* padre */
    printf("parent:%d ", x);
    wait(NULL);
    x++;
    printf("parent:%d ", x);
}
\end{lstlisting}

\subsubsection*{Análisis}

\texttt{fork()} crea un \textbf{nuevo proceso} con su propia copia de
la memoria. Después del \texttt{fork}, padre e hijo tienen cada uno
su propia variable \texttt{x} (originalmente 0 en ambos). Lo que uno
modifique \emph{no afecta} al otro.

\begin{itemize}[leftmargin=*,nosep]
    \item \textbf{Hijo}: hace \texttt{x += 1} (su $x$ vale 1),
        imprime \texttt{"child:1 "}, hace \texttt{exit(1)}.
    \item \textbf{Padre}: imprime \texttt{"parent:0 "} (su $x$ sigue
        en 0), entra a \texttt{wait(NULL)} hasta que el hijo termine,
        luego hace \texttt{x++} (su $x$ pasa a 1) e imprime
        \texttt{"parent:1 "}.
\end{itemize}

\textbf{¿Qué puede variar?} El orden entre el primer \texttt{printf}
del padre (\texttt{"parent:0 "}) y el del hijo (\texttt{"child:1 "})
depende del scheduling. El segundo \texttt{printf} del padre
(\texttt{"parent:1 "}) está \textbf{después} de \texttt{wait}, así que
siempre va al final.

\subsubsection*{Salidas posibles}

\begin{center}
\small
\begin{tabular}{@{}lp{8cm}@{}}
\toprule
Orden & Salida \\
\midrule
Padre imprime primero & \texttt{parent:0 child:1 parent:1\ } \\
Hijo imprime primero  & \texttt{child:1 parent:0 parent:1\ } \\
\bottomrule
\end{tabular}
\end{center}

\begin{idea}
\textbf{Lo que NO puede aparecer}: \texttt{"parent:1"} antes de
\texttt{"child:1"}, porque el segundo \texttt{printf} del padre está
después de \texttt{wait(NULL)}, que bloquea hasta que el hijo
finalice. Por construcción, \texttt{"parent:1"} es siempre lo último.
\end{idea}

\subsection*{Programa B (con threads)}

\begin{lstlisting}[style=cstyle]
int x = 0;

void f(void) {
    x++;
    printf("t1:%d ", x);
    exit_thread();
}

int main(void) {
    thread t1 = create_thread(f);
    printf("main:%d ", x);
    wait_thread(t1);
    printf("main:%d ", ++x);
}
\end{lstlisting}

\subsubsection*{Análisis}

A diferencia de \texttt{fork()}, \texttt{create\_thread()} crea un
\textbf{hilo dentro del mismo proceso}, que \textbf{comparte el
espacio de direcciones} con \texttt{main}. La variable \texttt{x} es
una sola: si el hilo \texttt{t1} hace \texttt{x++}, \texttt{main} la
ve modificada.

Esto introduce una \textbf{race condition} entre la modificación de
\texttt{x} en \texttt{t1} y la lectura de \texttt{x} en el primer
\texttt{printf} de \texttt{main}.

Análisis paso a paso:

\begin{itemize}[leftmargin=*]
    \item \texttt{t1} ejecuta: \texttt{x++} (queda en 1) $\to$
        \texttt{printf("t1:1 ")} $\to$ \texttt{exit\_thread}.
    \item \texttt{main} ejecuta: \texttt{printf("main:?")} (lee el
        $x$ actual, que puede ser 0 o 1 según si \texttt{t1} alcanzó
        a incrementar antes) $\to$ \texttt{wait\_thread} (espera fin
        de \texttt{t1}) $\to$ \texttt{++x} (sube $x$ a 2) $\to$
        \texttt{printf("main:2 ")}.
\end{itemize}

El \texttt{++x} del último \texttt{printf} se evalúa siempre
\emph{después} del \texttt{wait\_thread}, momento en el cual el
incremento de \texttt{t1} ya quedó hecho. Por eso el \textbf{último
valor impreso es siempre 2}.

\subsubsection*{Salidas posibles}

\begin{center}
\small
\begin{tabular}{@{}lll@{}}
\toprule
Interleaving & Valor de $x$ que lee main & Salida \\
\midrule
\texttt{t1} corre completo antes que main lea  & $1$ & \texttt{t1:1 main:1 main:2\ } \\
main lee antes que \texttt{t1} incremente      & $0$ & \texttt{main:0 t1:1 main:2\ } \\
\texttt{t1} incrementa, main lee, \texttt{t1} imprime & $1$ & \texttt{main:1 t1:1 main:2\ } \\
\bottomrule
\end{tabular}
\end{center}

\begin{importante}
La diferencia esencial entre los dos programas es la naturaleza de
\texttt{x}:

\begin{itemize}[leftmargin=*,nosep]
    \item En \textbf{programa A}, \texttt{fork} duplica el espacio de
        direcciones. Padre e hijo tienen variables \texttt{x}
        \emph{independientes}. No hay race condition.
    \item En \textbf{programa B}, los hilos \emph{comparten} el
        espacio de direcciones. La misma variable \texttt{x} es
        accedida concurrentemente. Hay race condition real: el valor
        que lee \texttt{main} depende del orden de ejecución.
\end{itemize}

Esta es exactamente la diferencia conceptual entre procesos y
threads.
\end{importante}

% =====================================================================
\section{Ejercicio 4 --- Análisis de deadlock (2 pts)}
% =====================================================================

\textbf{Enunciado.} Dada la siguiente secuencia de operaciones en un
sistema concurrente:

\begin{lstlisting}[style=textstyle]
1: T1: requiere y adquiere R1
2: T2: requiere R1
3: T1: libera R1
4: T3: requiere y adquiere R2
5: T2: adquiere R1
6: T1: requiere R2
7: T2: requiere R2
8: T3: requiere R1
9: ...
\end{lstlisting}

\textbf{(a)} Determinar si en el paso 8 se produce \emph{deadlock}.
Mostrar el grafo de dependencias.

\textbf{(b)} Si el administrador finaliza T2, ¿pueden las demás tareas
progresar?

\subsection*{Análisis: estado al finalizar el paso 8}

Reconstruyamos quién posee qué y quién espera qué después de cada
paso:

\begin{center}
\small
\begin{tabular}{@{}cllll@{}}
\toprule
Paso & Operación & R1 & R2 & Bloqueadas \\
\midrule
1 & T1 adquiere R1                 & T1 & ---  & --- \\
2 & T2 requiere R1 (T1 lo tiene)   & T1 & ---  & T2 \\
3 & T1 libera R1                   & libre & --- & T2 (despierta) \\
4 & T3 adquiere R2                 & libre & T3 & --- \\
5 & T2 adquiere R1                 & T2 & T3 & --- \\
6 & T1 requiere R2 (T3 lo tiene)   & T2 & T3 & T1 \\
7 & T2 requiere R2 (T3 lo tiene)   & T2 & T3 & T1, T2 \\
8 & T3 requiere R1 (T2 lo tiene)   & T2 & T3 & T1, T2, T3 \\
\bottomrule
\end{tabular}
\end{center}

\textbf{Configuración tras el paso 8:}

\begin{itemize}[leftmargin=*,nosep]
    \item R1 está asignado a T2.
    \item R2 está asignado a T3.
    \item T1 espera R2 (lo tiene T3).
    \item T2 espera R2 (lo tiene T3).
    \item T3 espera R1 (lo tiene T2).
\end{itemize}

\subsection*{Grafo de asignación de recursos}

Convención clásica: arista $R \to P$ si el recurso está asignado al
proceso; arista $P \to R$ si el proceso espera el recurso.

\begin{center}
\begin{tikzpicture}[
    node distance=2cm,
    proc/.style={circle,draw,minimum size=0.95cm,font=\bfseries\small},
    res/.style={rectangle,draw,minimum size=0.95cm,font=\bfseries\small},
    every edge/.style={draw,-{Stealth[length=2.5mm]},thick}
]
\node[proc]                            (T1) at (-3.5, 2.0) {$T_1$};
\node[proc]                            (T2) at (   0, 2.0) {$T_2$};
\node[proc]                            (T3) at ( 3.5, 2.0) {$T_3$};

\node[res]                             (R1) at ( -1.5, -1.0) {$R_1$};
\node[res]                             (R2) at (  1.5, -1.0) {$R_2$};

% asignaciones (R -> P)
\draw[->,blue!60!black,thick] (R1) -- node[below right,font=\scriptsize]{tiene} (T2);
\draw[->,blue!60!black,thick] (R2) -- node[below right,font=\scriptsize]{tiene} (T3);

% pedidos (P -> R)
\draw[->,red!70!black,thick,dashed] (T1) -- node[above left,font=\scriptsize]{espera} (R2);
\draw[->,red!70!black,thick,dashed] (T2) -- node[left,font=\scriptsize]{espera} (R2);
\draw[->,red!70!black,thick,dashed] (T3) -- node[above right,font=\scriptsize]{espera} (R1);
\end{tikzpicture}
\end{center}

\subsection*{4.a --- ¿Hay deadlock?}

\textbf{Sí}. Sigamos las aristas a partir de T2:

\quad T2 $\to$ R2 (espera) $\to$ T3 (tiene) $\to$ R1 (espera) $\to$ T2 (tiene) $\to$ \ldots

Hay un \textbf{ciclo cerrado} $T_2 \to R_2 \to T_3 \to R_1 \to T_2$.
Como los recursos son de instancia única, un ciclo basta para
garantizar deadlock. Las dos tareas atrapadas en el ciclo son
\textbf{T2 y T3}.

\textbf{¿Y T1?} T1 espera R2. R2 está retenido por T3, que a su vez
está en deadlock. T1 \textbf{no forma parte del ciclo} pero está
\emph{bloqueada como víctima colateral}: nunca obtendrá R2 porque T3
nunca lo va a liberar. T1 es lo que se llama una \emph{víctima
indirecta} del deadlock.

\begin{respuesta}
\textbf{(a)} Hay deadlock entre \textbf{T2 y T3} (ciclo $T_2 \to R_2
\to T_3 \to R_1 \to T_2$). \textbf{T1} también está bloqueada pero
como consecuencia secundaria del deadlock, no como integrante del
ciclo.
\end{respuesta}

\subsection*{4.b --- Si el admin finaliza T2}

Si el administrador finaliza T2 (por ejemplo, con \texttt{kill -9}):

\begin{enumerate}[leftmargin=*,nosep]
    \item T2 muere y libera R1 (todo recurso retenido por un proceso
        se libera cuando éste termina).
    \item T3 estaba esperando R1 $\to$ ahora puede adquirirlo. T3
        ahora tiene R1 y R2.
    \item T3 progresa, termina su sección crítica y libera R1 y R2.
    \item T1 estaba esperando R2 $\to$ al liberarse, T1 lo adquiere y
        continúa hasta terminar.
\end{enumerate}

\begin{respuesta}
\textbf{(b)} Sí, las demás tareas pueden progresar. Finalizar T2
rompe el ciclo del deadlock: T3 obtiene R1 y termina su trabajo, lo
que a su vez permite que T1 obtenga R2. Ésta es la técnica clásica de
\textbf{recuperación de deadlock} matando un proceso del ciclo.
\end{respuesta}

% =====================================================================
\section{Ejercicio 5 --- Planificación con prioridades y \emph{preemption}}
% =====================================================================

\textbf{Enunciado.} Dadas las siguientes tareas, todas creadas en
$t=0$:

\begin{center}
\small
\begin{tabular}{@{}lcl@{}}
\toprule
Tarea & Prioridad & Comportamiento \\
\midrule
\texttt{idle} & 0 & \texttt{CPU:1 CPU:1 ...} \\
\texttt{P1}   & 1 & \texttt{CPU:2.5 EXIT} \\
\texttt{P2}   & 2 & \texttt{CPU:1 WAIT:1 CPU:1 EXIT} \\
\bottomrule
\end{tabular}
\end{center}

Scheduler basado en prioridades con \emph{preemption} y
\texttt{quantum}=2.

\textbf{Convención de prioridad.} Bajo la asignación dada, \texttt{idle}
tiene prioridad 0 (la mínima), \texttt{P1} prioridad 1 y \texttt{P2}
prioridad 2 (la máxima). \emph{Mayor número $\to$ mayor prioridad.}

\subsection*{5.a --- Orden de la cola RUNNABLE tras el arribo}

Al instante $t=0$ las tres tareas entran simultáneamente al sistema.
La cola \texttt{RUNNABLE} (o \texttt{READY}) se ordena por prioridad
\textbf{descendente}:

\begin{respuesta}
\textbf{(a)} Orden de la cola \texttt{READY} en $t = 0$:
\begin{center}\Large \texttt{[P2, P1, idle]}\end{center}
P2 con prioridad 2 toma la CPU primero; P1 (prio 1) queda en segundo
lugar; \texttt{idle} (prio 0) al final.
\end{respuesta}

\subsection*{5.b --- ¿P1 termina antes que P2?}

Simulamos paso a paso. Usaremos la notación \texttt{[t$_i$,t$_f$)} =
intervalo en que cada tarea corre.

\paragraph{Paso 1: \texttt{[0,1)} --- P2 ejecuta su primer
\texttt{CPU:1}.}
P2 (prio 2) es la de mayor prioridad. Corre 1 unidad. Al llegar a
$t=1$, P2 termina su primera ráfaga y entra a \texttt{WAIT:1} (se
duerme por 1 unidad).

\paragraph{Paso 2: \texttt{[1,2)} --- P1 toma la CPU.}
Con P2 dormida, la \texttt{READY} queda \texttt{[P1, idle]}. P1 (prio
1) supera a \texttt{idle} (prio 0). P1 ejecuta 1 unidad de su
\texttt{CPU:2.5}, le quedan $2.5 - 1 = 1.5$.

\paragraph{Paso 3: \texttt{[2,2)} --- P2 despierta y \emph{preempta} a
P1.}
En $t=2$, P2 termina su \texttt{WAIT:1} y vuelve a \texttt{READY}. Como
P2 tiene mayor prioridad que P1, \textbf{preempta} a P1. P1 vuelve a
\texttt{READY} con 1.5 unidades pendientes.

\paragraph{Paso 4: \texttt{[2,3)} --- P2 ejecuta su segundo
\texttt{CPU:1}.}
P2 corre 1 unidad más y al llegar a $t=3$ ejecuta \texttt{EXIT}. P2
termina en $t=3$.

\paragraph{Paso 5: \texttt{[3,4.5)} --- P1 retoma y termina.}
Con P2 fuera, la \texttt{READY} es \texttt{[P1, idle]}. P1 toma la
CPU. Le faltan 1.5 unidades. El \emph{quantum} es 2, así que no se
agotará por tiempo. P1 corre 1.5 unidades y en $t=4.5$ hace
\texttt{EXIT}.

\subsection*{Diagrama de Gantt}

\begin{center}
\begin{tikzpicture}[xscale=1.6, yscale=0.7,
    every node/.style={font=\ttfamily\small}]
\draw[fill=c2!50] (0,0) rectangle (1,1) node[midway,white,font=\bfseries]{P2};
\draw[fill=c1!50] (1,0) rectangle (2,1) node[midway,white,font=\bfseries]{P1};
\draw[fill=c2!50] (2,0) rectangle (3,1) node[midway,white,font=\bfseries]{P2};
\draw[fill=c1!50] (3,0) rectangle (4.5,1) node[midway,white,font=\bfseries]{P1};
\draw[fill=c3!30] (4.5,0) rectangle (6,1) node[midway,font=\bfseries]{idle\ldots};

% marcas de tiempo
\foreach \x in {0,1,2,3,4.5,6}
    \draw (\x,0) -- (\x,-0.2) node[below,font=\scriptsize]{\x};

% anotaciones
\node[font=\scriptsize,below=0.5cm] at (0.5,-0.1) {CPU:1};
\node[font=\scriptsize,below=0.5cm] at (1.5,-0.1) {1 de 2.5};
\node[font=\scriptsize,below=0.5cm] at (2.5,-0.1) {CPU:1};
\node[font=\scriptsize,below=0.5cm] at (3.75,-0.1) {1.5 restantes};

% estados de P2: WAIT
\draw[thick,dashed,c2!70!black] (1,1.4) -- (2,1.4);
\node[font=\scriptsize,above] at (1.5,1.4) {P2 en WAIT};
\end{tikzpicture}
\end{center}

\subsection*{Métricas (turnaround)}

\begin{center}
\small
\begin{tabular}{@{}lccccc@{}}
\toprule
Tarea & Llegada & Fin & Turnaround & CPU usada & Espera \\
\midrule
P2   & 0 & 3   & 3   & 2  & 0 \\
P1   & 0 & 4.5 & 4.5 & 2.5 & 1 \\
idle & 0 & --- (nunca termina) & --- & --- & 4.5 \\
\bottomrule
\end{tabular}
\end{center}

\begin{respuesta}
\textbf{(b)} \textbf{NO}. P1 \emph{no} termina antes que P2:

\begin{itemize}[leftmargin=*,nosep]
    \item P2 termina en $t = 3$.
    \item P1 termina en $t = 4.5$.
\end{itemize}

\textbf{Razón}: aunque P2 tiene una ráfaga de CPU más corta (2
unidades en total vs.\ 2.5 de P1), su \textbf{prioridad superior}
hace que siempre que esté \texttt{READY} preempte a P1. Cuando P2 va
a \texttt{WAIT}, P1 progresa, pero cuando P2 despierta vuelve a
quitarle la CPU.
\end{respuesta}

\subsection*{5.c --- ¿Cuándo se planifica la tarea \texttt{idle}?}

La tarea \texttt{idle} (prio 0) solo es planificada cuando \emph{no
hay ninguna otra tarea} en \texttt{READY}:

\begin{itemize}[leftmargin=*,nosep]
    \item De $t=0$ a $t=4.5$ siempre hay alguna tarea de mayor
        prioridad (\texttt{P1} o \texttt{P2}) en estado
        \texttt{READY} o \texttt{RUNNING}. Aunque P2 esté en
        \texttt{WAIT} durante \texttt{[1,2)}, P1 ocupa la CPU en ese
        intervalo, así que \texttt{idle} no corre.
    \item En $t=3$ P2 termina. P1 sigue activa.
    \item En $t=4.5$ P1 termina. La cola \texttt{READY} queda solo
        con \texttt{idle}. Recién entonces \texttt{idle} toma la CPU.
\end{itemize}

\begin{respuesta}
\textbf{(c)} \texttt{idle} se planifica desde $t = 4.5$ en adelante,
una vez que \texttt{P1} y \texttt{P2} han terminado y no hay otra
tarea \texttt{READY}. Es la función clásica de \texttt{idle}:
ocupar la CPU cuando no hay trabajo útil para hacer.
\end{respuesta}

% =====================================================================
\section*{Resumen del parcial}
\addcontentsline{toc}{section}{Resumen del parcial}
% =====================================================================

\begin{center}
\small
\begin{tabular}{@{}clp{8cm}@{}}
\toprule
\# & Tema & Respuesta sintética \\
\midrule
1 & V/F sobre SO    & a: V, b: F, c: F, d: V, e: F, f: V, g: F, h: F \\
2 & Linker estático & a: F (dir.\ \texttt{f} = \texttt{0x00FA}); b: F (call externo); c: V (\texttt{0x0184}) \\
3 & \texttt{fork} vs threads & A: dos órdenes posibles; B: tres interleavings, todos terminan con \texttt{main:2} \\
4 & Deadlock        & (a) Sí, ciclo $T_2 \to R_2 \to T_3 \to R_1 \to T_2$; T1 víctima colateral. (b) Sí, finalizar T2 desbloquea a T3 y T1 \\
5 & Scheduling con prioridades & (a) \texttt{[P2,P1,idle]}; (b) NO: P2 termina en $t=3$, P1 en $t=4.5$; (c) \texttt{idle} a partir de $t=4.5$ \\
\bottomrule
\end{tabular}
\end{center}

\end{document}
